% Part: first-order-logic
% Chapter: syntax-and-semantics
% Section: formation-sequences

\documentclass[../../../include/open-logic-section]{subfiles}

\begin{document}

\olfileid{pl}{syn}{fseq}

\olsection{متتاليات التكوين}

من أحسن طرق البحث أن تُعرّف !!{formula}s استقرائيًا، ثم تُثبت خواص
!!{formula}s بالاستقراء. ومع حسن هذه الطريقة وقلة مؤونتها، قد ينفع أن نتابع بناء
!!{formula}s من مبادئه خطوة خطوة. ولأجل هذا نأخذ الآن بمفهوم
\emph{متتالية التكوين}.

\begin{defn}[متتاليات التكوين للصيغ]
\ollabel{defn:fseq-frm}
تكون المتتالية المنتهية $\tuple{!A_{\OLMathNumeral{0}},\dotsc,!A_{\OLId{latin-ordinary}{n}}}$ المكوّنة من سلاسل
رموز من اللغة~$\Lang L_{\OLMathNumeral{0}}$ \emph{متتالية تكوين} لـ~$!A$ إذا كان
$!A \ident !A_{\OLId{latin-ordinary}{n}}$، وكان، لكل ${\OLId{latin-ordinary}{i}} \leq {\OLId{latin-ordinary}{n}}$، إما أن تكون $!A_{\OLId{latin-ordinary}{i}}$ صيغة ذرية،
وإما أن يوجد ${\OLId{latin-ordinary}{j}},{\OLId{latin-ordinary}{k}} < {\OLId{latin-ordinary}{i}}$ بحيث تتحقق إحدى الحالات الآتية:
\begin{enumerate}
    \tagitem{prvNot}{$!A_{\OLId{latin-ordinary}{i}} \ident \lnot !A_{\OLId{latin-ordinary}{j}}$.}{}%
    \tagitem{prvAnd}{$!A_{\OLId{latin-ordinary}{i}} \ident (!A_{\OLId{latin-ordinary}{j}} \land !A_{\OLId{latin-ordinary}{k}})$.}{}%
    \tagitem{prvOr}{$!A_{\OLId{latin-ordinary}{i}} \ident (!A_{\OLId{latin-ordinary}{j}} \lor !A_{\OLId{latin-ordinary}{k}})$.}{}%
    \tagitem{prvIf}{$!A_{\OLId{latin-ordinary}{i}} \ident (!A_{\OLId{latin-ordinary}{j}} \lif !A_{\OLId{latin-ordinary}{k}})$.}{}%
    \tagitem{prvIff}{$!A_{\OLId{latin-ordinary}{i}} \ident (!A_{\OLId{latin-ordinary}{j}} \liff !A_{\OLId{latin-ordinary}{k}})$.}{}%
\end{enumerate}
\end{defn}

\begin{ex}
\[
\tuple{
    \Obj p_{\OLMathNumeral{0}},
    \Obj p_{\OLMathNumeral{1}},
    (\Obj p_{\OLMathNumeral{1}} \land \Obj p_{\OLMathNumeral{0}}),
    \lnot (\Obj p_{\OLMathNumeral{1}} \land \Obj p_{\OLMathNumeral{0}})
}
\]
هي متتالية تكوين لـ
$\lnot (\Obj p_{\OLMathNumeral{1}} \land \Obj p_{\OLMathNumeral{0}})$، وكذلك
\[
\tuple{
    \Obj p_{\OLMathNumeral{0}},
    \Obj p_{\OLMathNumeral{1}},
    \Obj p_{\OLMathNumeral{0}},
    (\Obj p_{\OLMathNumeral{1}} \land \Obj p_{\OLMathNumeral{0}}),
    (\Obj p_{\OLMathNumeral{0}} \lif \Obj p_{\OLMathNumeral{1}}),
    \lnot (\Obj p_{\OLMathNumeral{1}} \land \Obj p_{\OLMathNumeral{0}})
}.
\]
%
وكما يتبين من المثال الثاني، قد تحتوي متتاليات التكوين على «حشو»:
أي صيغ زائدة أو لا تسهم في البناء.
\end{ex}

\begin{prop}
\ollabel{prop:formed}
لكل !!{formula}~$!A$ في~$\Frm[{\OLId{latin-looped}{L}}_{\OLMathNumeral{0}}]$ متتالية تكوين.
\end{prop}

\begin{proof}
افترض أن $!A$ ذرية. عندئذ تكون المتتالية~$\tuple{!A}$ متتالية
تكوين لـ~$!A$.
%
وافترض الآن أن لـ~$!B$ و~$!C$ متتاليتا تكوين
$\tuple{!B_{\OLMathNumeral{0}},\dotsc,!B_{\OLId{latin-ordinary}{n}}}$ و~$\tuple{!C_{\OLMathNumeral{0}},\dotsc,!C_{\OLId{latin-ordinary}{m}}}$،
على الترتيب.
%
\begin{enumerate}
  \tagitem{prvNot}{إذا كان $!A \ident \lnot !B$،
      فتكون $\tuple{!B_{\OLMathNumeral{0}},\dotsc,!B_{\OLId{latin-ordinary}{n}},\lnot !B_{\OLId{latin-ordinary}{n}}}$
      متتالية تكوين لـ~$!A$.}{}
  \tagitem{prvAnd}{إذا كان $!A \ident (!B \land !C)$،
      فتكون $\tuple{!B_{\OLMathNumeral{0}},\dotsc,!B_{\OLId{latin-ordinary}{n}},!C_{\OLMathNumeral{0}},\dotsc,!C_{\OLId{latin-ordinary}{m}},(!B_{\OLId{latin-ordinary}{n}} \land !C_{\OLId{latin-ordinary}{m}})}$
      متتالية تكوين لـ~$!A$.}{}
  \tagitem{prvOr}{إذا كان $!A \ident (!B \lor !C)$،
      فتكون $\tuple{!B_{\OLMathNumeral{0}},\dotsc,!B_{\OLId{latin-ordinary}{n}},!C_{\OLMathNumeral{0}},\dotsc,!C_{\OLId{latin-ordinary}{m}},(!B_{\OLId{latin-ordinary}{n}} \lor !C_{\OLId{latin-ordinary}{m}})}$
      متتالية تكوين لـ~$!A$.}{}
  \tagitem{prvIf}{إذا كان $!A \ident (!B \lif !C)$،
      فتكون $\tuple{!B_{\OLMathNumeral{0}},\dotsc,!B_{\OLId{latin-ordinary}{n}},!C_{\OLMathNumeral{0}},\dotsc,!C_{\OLId{latin-ordinary}{m}},(!B_{\OLId{latin-ordinary}{n}} \lif !C_{\OLId{latin-ordinary}{m}})}$
      متتالية تكوين لـ~$!A$.}{}
  \tagitem{prvIff}{إذا كان $!A \ident (!B \liff !C)$،
      فتكون $\tuple{!B_{\OLMathNumeral{0}},\dotsc,!B_{\OLId{latin-ordinary}{n}},!C_{\OLMathNumeral{0}},\dotsc,!C_{\OLId{latin-ordinary}{m}},(!B_{\OLId{latin-ordinary}{n}} \liff !C_{\OLId{latin-ordinary}{m}})}$
      متتالية تكوين لـ~$!A$.}{}
\end{enumerate}
وبحسب مبدأ الاستقراء على !!{formula}s، لكل
!!{formula} متتالية تكوين.
\end{proof}

وأما العكس فيصح أيضًا، وإثباته يبيّن أن طريقتي تعريف الصيغ متكافئتان، لا تختلف نتائجهما. ومن ثمرته أن نثبت مبرهنات الصيغ بالاستقراء العادي على أطوال متتاليات تكوينها.

\begin{lem}
\ollabel{lem:fseq-init}
افترض أن $\tuple{!A_{\OLMathNumeral{0}},\dotsc,!A_{\OLId{latin-ordinary}{n}}}$ متتالية تكوين لـ~$!A_{\OLId{latin-ordinary}{n}}$، وأن
${\OLId{latin-ordinary}{k}} \leq {\OLId{latin-ordinary}{n}}$. عندئذ تكون $\tuple{!A_{\OLMathNumeral{0}},\dotsc,!A_{\OLId{latin-ordinary}{k}}}$ متتالية تكوين
لـ~$!A_{\OLId{latin-ordinary}{k}}$.
\end{lem}

\begin{proof}
تمرين.
\end{proof}

\begin{thm}
\ollabel{thm:fseq-frm-equiv}
$\Frm[{\OLId{latin-looped}{L}}_{\OLMathNumeral{0}}]$ هي مجموعة جميع سلاسل الرموز في اللغة~$\Lang L_{\OLMathNumeral{0}}$ التي
لها متتالية تكوين.
\end{thm}

\begin{proof}
اجعل ${\OLId{latin-looped}{F}}$ مجموعة سلاسل رموز اللغة~$\Lang L_{\OLMathNumeral{0}}$ التي لها متتالية تكوين.
قد ثبت في \olref[pl][syn][fseq]{prop:formed} أن
$\Frm[{\OLId{latin-looped}{L}}_{\OLMathNumeral{0}}] \subseteq {\OLId{latin-looped}{F}}$؛ فليس علينا إلا إثبات الاحتواء الآخر.

لتكن لـ~$!A$ متتالية تكوين $\tuple{!A_{\OLMathNumeral{0}},\dotsc,!A_{\OLId{latin-ordinary}{n}}}$.
والمطلوب $!A \in \Frm[{\OLId{latin-looped}{L}}_{\OLMathNumeral{0}}]$؛ فنستقرئ استقراء قويًا على~${\OLId{latin-ordinary}{n}}$.
نفرض أن كل سلسلة ذات متتالية تكوين رقم حدّها الأخير
${\OLId{latin-ordinary}{m}} < {\OLId{latin-ordinary}{n}}$ تقع في~$\Frm[{\OLId{latin-looped}{L}}_{\OLMathNumeral{0}}]$. وتعريف متتالية التكوين يقسم الأمر: إما أن $!A_{\OLId{latin-ordinary}{n}}$ ذرية، وإما أن يوجد ${\OLId{latin-ordinary}{j}},{\OLId{latin-ordinary}{k}} < {\OLId{latin-ordinary}{n}}$ وتتحقق إحدى الصور الآتية:
\begin{enumerate}
    \tagitem{prvNot}{$!A_{\OLId{latin-ordinary}{n}} \ident \lnot !A_{\OLId{latin-ordinary}{j}}$.}{}%
    \tagitem{prvAnd}{$!A_{\OLId{latin-ordinary}{n}} \ident (!A_{\OLId{latin-ordinary}{j}} \land !A_{\OLId{latin-ordinary}{k}})$.}{}%
    \tagitem{prvOr}{$!A_{\OLId{latin-ordinary}{n}} \ident (!A_{\OLId{latin-ordinary}{j}} \lor !A_{\OLId{latin-ordinary}{k}})$.}{}%
    \tagitem{prvIf}{$!A_{\OLId{latin-ordinary}{n}} \ident (!A_{\OLId{latin-ordinary}{j}} \lif !A_{\OLId{latin-ordinary}{k}})$.}{}%
    \tagitem{prvIff}{$!A_{\OLId{latin-ordinary}{n}} \ident (!A_{\OLId{latin-ordinary}{j}} \liff !A_{\OLId{latin-ordinary}{k}})$.}{}%
\end{enumerate}
فإن كانت $!A_{\OLId{latin-ordinary}{n}}$ ذرية، فقد حصل
$!A_{\OLId{latin-ordinary}{n}} \in \Frm[{\OLId{latin-looped}{L}}_{\OLMathNumeral{0}}]$. وإلا فلنأخذ حالة $!A \ident
(!A_{\OLId{latin-ordinary}{j}} \land !A_{\OLId{latin-ordinary}{k}})$. تقضي \olref[pl][syn][fseq]{lem:fseq-init} بأن
$\tuple{!A_{\OLMathNumeral{0}},\dotsc,!A_{\OLId{latin-ordinary}{j}}}$ و~$\tuple{!A_{\OLMathNumeral{0}},\dotsc,!A_{\OLId{latin-ordinary}{k}}}$ متتاليتا تكوين لـ~$!A_{\OLId{latin-ordinary}{j}}$ و~$!A_{\OLId{latin-ordinary}{k}}$ على الترتيب. وهما مقطعان ابتدائيان حقيقيان من متتالية~$!A$؛ فرقم الحد الأخير في كل منهما أقل من~${\OLId{latin-ordinary}{n}}$. وبفرض الاستقراء يقع~$!A_{\OLId{latin-ordinary}{j}}$ و~$!A_{\OLId{latin-ordinary}{k}}$ في
$\Frm[{\OLId{latin-looped}{L}}_{\OLMathNumeral{0}}]$. ومن تعريف !!a{formula} يلزم أن
$(!A_{\OLId{latin-ordinary}{j}} \land !A_{\OLId{latin-ordinary}{k}})$ كذلك. وسائر الحالات تُحمل على هذا المسلك.
\end{proof}

\end{document}
